Перед формулировкой и доказательством основных теорем нам понадобятся
следующие предварительные результаты.
\subsection{Система аксиом NBG.}
Настоящая работа проводится в рамках системы аксиом теории множеств
фон Неймана -- Бернайса -- Гёделя. В данной системе аксиом все
объекты называются классами. Нам понадобится следующее определение.
\begin{definition}[Собственный класс]
  Класс \( A \) называется \emph{собственным}, если он не является
  элементом никакого другого класса.
\end{definition}
Все классы, не являющиеся собственными, являются множествами.
Приведем некоторые примеры:
\begin{enumerate}
  \item Класс всех множеств является собственным.
  \item Класс всех кардиналов является собственным.
\end{enumerate}
\begin{lemma}
  Класс всех метрических пространств является собственным.
\end{lemma}
\begin{proof}
  На любом множестве \( X \) можно задать дискретную метрику \( d \),
  \( d(x,y) = 1 \) для всех \( x,y \in X \), \( x,y\neq 0 \) и \(
  d(x,x) = 0 \) для всех \( x\in X \). Получаем, что все множества
  лежат в классе всех метрических пространств, а так как класс всех
  множеств собственный, то и класс всех метрических пространств также
  является собственным.
\end{proof}
В дальнейшем будет показано, что основные объекты изучения данной
работы, облака, являются собственными классами.
\subsection{Основные определения.}

\begin{definition}[Симплекс]
  Симплексом мощности \( \alpha \) называется метрическое
  пространство \( \Delta _
  \alpha  \) мощности \( \alpha  \) с дискретной метрикой. В
  частности \( \Delta _1 \) будет обозначать одноточечное пространство.
\end{definition}
Пусть $X$ и $Y$ ---
метрические пространства. Расстояние между точками метрического
пространства \( x, x^\prime \) будем обозначать \( | xx^\prime | \).
Между метрическими пространствами можно задать расстояние, называемое
расстоянием Громова--Хаусдорфа.
\begin{definition}[Расстояние Хаусдорфа]
  Пусть \( X,Y \) --- метрические подпространства метрического
  пространства \( Z \). \emph{Расстоянием Хаусдорфа} между
  пространствами \( X \) и \( Y \) называется величина
  \[
    d_H(X,Y) = \max \left\{\sup_{x\in X}\inf_{y\in Y} | xy |,
    \sup_{y\in Y}\inf_{x\in X} | xy |\right\}.
  \]
\end{definition}
\begin{definition}[Расстояние Громова-Хаусдорфа]
  \emph{Реализацией} метрических пространств \( X, Y \) называется
  тройка метрических пространств \( X', Y', Z \) таких, что \( X'
  \subset Z \) , \( Y' \subset Z \), \( X \cong X' \), \( Y \cong Y' \).

  \emph{Расстоянием Громова-Хаусдорфа}  $d_{GH}(X,Y)$ между метрическими
  пространствами $X, Y$ является точная нижняя  грань чисел $r$
  таких, что существует реализация $(X',Y',Z)$ и $d_H(X', Y')  \le
  r$, где $d_H$ --- расстояние Хаусдорфа.
\end{definition}
Мы также будем пользоваться другим эквивалентным определением
расстояния Громова-Хаусдорфа \cite{Lectures}.
\begin{definition}[Соответствие]
  Пусть $X$, $Y$ --- метрические пространства. \emph{Соответствием}
  $R$ между этими пространствами
  называется сюръективное многозначное отображение между ними.
  Множество всех соответствий между $X$ и $Y$ обозначается
  $\mathcal{R}(X,Y)$. Также будем отождествлять соответствие и его
  график.

  \emph{Искажением} соответствия \( R \) называется величина
  \[
    \dis{R} = \sup{\Bigl\{ \big| |xx'| - |yy'| \big| : (x, y), (x',
    y') \in R\Bigr\}}.
  \]
\end{definition}
\begin{definition}[Расстояние Громова-Хаусдорфа]
  \emph{Расстояние Громова-Хаусдорфа} между метрическими
  пространствами будем называть величину
  \[
    d_{GH}(X,Y) = \frac{1}{2}\inf \bigl\{\dis{R} : R \in
    \mathcal{R}(X,Y)\bigr\}.
  \]
\end{definition}
\begin{lemma}
  \label{lem:GHproperties}
  Расстояние Громова-Хаусдорфа обладает следующими свойствами.
  \begin{enumerate}
    \item \( d_{GH}(X,Y) \le \frac{1}{2}\max\{\diam X, \diam Y\} \).
    \item \( d_{GH}(X,Y) \ge \frac{1}{2} | \diam X - \diam Y |\).
      \label{lem:GHproperties:2}
  \end{enumerate}
\end{lemma}
Расстояние Громова-Хаусдорфа на собственном классе всех метрических
пространств, рассматриваемых с точностью до изометрии является
обобщенной псевдометрикой. Действительно, оно неотрицательно, \(
d_{GH}(X,X) = 0 \) для всякого метрического пространства \( X \), и
для него выполняется неравенство треугольника.
\begin{example} Следующие примеры показывают, почему расстояние
  Громова-Хаусдорфа --- обобщенная псевдометрика.
  \begin{enumerate}
    \item \( d_{GH}\big([0,1],(0,1)\big) =0\) так как естественное
      вложение \( (0,1) \hookrightarrow [0,1] \) имеет расстояние
      Хаусдорфа \( 0 \). При этом пространства не изометричны так как
      только \( [0,1] \) является компактом.
    \item \( d_{GH}\left([0,1], \R\right)  = \infty\) по свойству
      \ref{lem:GHproperties:2} из леммы \ref{lem:GHproperties}.
  \end{enumerate}
\end{example}
Далее будем рассматривать метрические пространства с точностью до
нулевого расстояния между ними. Получившийся класс обозначим \(
\mathcal{GH}_0 \). На нем расстояние Громова-Хаусдорфа становится
обобщенной метрикой.

\begin{definition}[\cite{TuzhBog1}] В классе $\mathcal{GH}_{0}$
  рассмотрим следующее
  отношение: $X \thicksim Y \Leftrightarrow d_{GH}(X, Y) < \infty$. Нетрудно
  убедиться, что оно будет отношением эквивалентности. Классы этой
  эквивалентности
  называются \emph{облаками}. Облако, в котором лежит метрическое
  пространство $X$
  будем обозначать $[X]$.
\end{definition}

Для любого метрического пространства $X$ определена операция умножения его
на положительное вещественное число $\lambda\colon X\mapsto \lambda X$, а именно
$(X, \rho) \mapsto (X, \lambda \rho)$, расстояние между любыми точками
пространства изменяется в $\lambda$ раз.
\begin{remark} Пусть метрические
  пространства $X$, $Y$ лежат в одном облаке. Тогда
  $d_{GH}(\lambda X, \lambda Y) = \lambda d_{GH}(X,Y) < \infty$, т.е.
  пространства
  $\lambda X$, $\lambda Y$ также будут лежать в одном облаке.
  \label{remOneCloud}
\end{remark}
\begin{definition}Определим операцию умножения облака $[X]$ на
  положительное вещественное
  число $\lambda$ как отображение, переводящее все пространства $Y \in [X]$ в
  пространства $\lambda Y$. По замечанию $\ref{remOneCloud}$ все
  полученные пространства будут
  лежать в облаке $[\lambda X]$.
\end{definition}

При таком отображении облако может как
измениться, так и перейти в себя. Для последнего случая вводится специальное
определение.
\begin{definition}[\cite{TuzhBog2}]
  Стационарной группой $\St\bigl([X]\bigr)$ облака $[X]$
  называется подмножество $\mathbb{R}_+$ такое, что для всех
  $\lambda \in \St\bigl([X]\bigr)$, $[X] = [\lambda X]$. Полученное подмножество
  действительно будет подгруппой в $\mathbb{R}_+$. Тривиальной
  будем называть стационарную группу равную $\{1\}$.
\end{definition}

\begin{example}
  Приведем несколько примеров облаков и их стационарных групп.
  \begin{enumerate}
    \item Пусть $\Delta_1$ --- одноточечное метрическое
      пространство.  Тогда\\ $\St\bigl([\Delta_1]\bigr) = \mathbb{R}_+$.
    \item $\St\bigl([\mathbb{R}]\bigr) = \mathbb{R}_+ $.
    \item Предположим, что функция $\varphi(n)$ удовлетворяет соотношению
      $\lim\limits_{n \rightarrow \infty } \varphi(n + 1) - \varphi(n) =
      +\infty$. Для $q > 1$ зададим пространство $X_q =
      \left\{q^{\varphi(n)}: n \in \mathbb{N}\right\} \subset \R$ с метрикой,
      индуцированной из \( \R \).
      Тогда выполняется
      $\St\bigl(\left[X_q\right]\bigr) = \{1\}$\cite{TuzhBog1}.
    \item Для натурального $p$ зададим пространство $X_p = \left\{p^n:n
      \in \mathbb{Z}\right\} \subset \R$ с метрикой, индуцированной из \( \R \).
      Для любого простого $p$ выполняется $\St\bigl(\left[X_p\right]\bigr) =
      \left\{p^n:n \in \mathbb{Z}\right\}$\cite{BogBog1}.
  \end{enumerate}
\end{example}
\begin{lemma}[\cite{TuzhBog2}]
  В каждом облаке с нетривиальной стационарной группой существует единственное
  пространство $X$ такое, что для любого $\lambda$ из стационарной группы
  выполняется $X = \lambda X$.
  \label{centerLemma}
\end{lemma}
\begin{definition}
  Пространство из леммы \ref{centerLemma} будем называть \emph{центром} облака.
\end{definition}
\begin{remark} В
  облаке $[\Delta_1]$ для любого пространства $X$ выполняется:
  $$|\lambda X, \mu X| = |\lambda - \mu||X,\Delta_1|.$$
\end{remark}
\begin{remark}[Ультраметрическое неравенство] В облаке
  $[\Delta_{1}]$ для всех пространств $X_{1}, X_{2}$ выполняется неравенство:
  \[
    |X_{1},X_{2}| \le \max\big\{|X_{1}, \Delta_{1}|,|X_{2},\Delta_{1}|\big\}.
  \]
  \label{remUltraMetric}
\end{remark}
